\documentclass[aspectratio=169,11pt]{beamer}

\usetheme{default}
\setbeamertemplate{navigation symbols}{}
\setbeamertemplate{footline}[frame number]

\usepackage{amsmath}
\usepackage{amsfonts}
\usepackage{booktabs}
\usepackage{graphicx}
\graphicspath{{master/figures/}{figures/}}

\definecolor{kmblue}{RGB}{30,78,121}
\definecolor{kmgold}{RGB}{193,148,37}
\definecolor{kmslate}{RGB}{59,70,84}
\definecolor{kmlight}{RGB}{240,244,248}

\setbeamercolor{title}{fg=white,bg=kmblue}
\setbeamercolor{frametitle}{fg=white,bg=kmblue}
\setbeamercolor{structure}{fg=kmblue}
\setbeamercolor{block title}{fg=white,bg=kmblue}
\setbeamercolor{block body}{fg=black,bg=kmlight}
\setbeamercolor{normal text}{fg=black,bg=white}
\setbeamercolor{itemize item}{fg=kmgold}
\setbeamercolor{itemize subitem}{fg=kmgold}
\setbeamerfont{title}{series=\bfseries,size=\LARGE}
\setbeamerfont{frametitle}{series=\bfseries,size=\large}

\title{K-Means Clustering}
\subtitle{From intuition to objective functions and model selection}
\author{}
% \date{March 27, 2026}

\newcommand{\wcss}{\mathrm{WCSS}}

\begin{document}

\begin{frame}[plain]
  \titlepage
  \vspace{-0.5em}
  \begin{center}
    % \includegraphics[
    %   width=0.48\linewidth,
    %   height=0.32\textheight,
    %   keepaspectratio
    % ]{intro_clustering_overview.jpg}
  \end{center}
\end{frame}

\begin{frame}[t]{Why We Cluster}
  \begin{columns}[T,totalwidth=\textwidth]
    \column{0.38\textwidth}
      \begin{itemize}
        \item Clustering is \textbf{unsupervised}: we observe feature vectors, but no labels.
        \item The goal is to create groups with high within-cluster similarity and low between-cluster similarity.
        \item Common use cases include customer segmentation, document organization, image grouping, social network analysis, and genomics.
      \end{itemize}
      \vspace{0.35em}
      \[
        \text{No labels} \Longrightarrow \text{discover structure from geometry alone}
      \]
    \column{0.58\textwidth}
      \begin{center}
        \includegraphics[
          width=\linewidth,
          height=0.60\textheight,
          keepaspectratio
        ]{intro_clustering_overview.jpg}
      \end{center}
  \end{columns}
\end{frame}

\begin{frame}[t]{Clustering Shows Up in Many Domains}
  \begin{columns}[T,totalwidth=\textwidth]
    \column{0.54\textwidth}
      \begin{center}
        \textbf{Documents and networks}

        \vspace{0.3em}
        \includegraphics[
          width=\linewidth,
          height=0.56\textheight,
          keepaspectratio
        ]{domain_documents_networks.jpg}
      \end{center}
    \column{0.42\textwidth}
      \begin{center}
        \textbf{Images and visual categories}

        \vspace{0.3em}
        \includegraphics[
          width=\linewidth,
          height=0.56\textheight,
          keepaspectratio
        ]{domain_visual_categories.jpg}
      \end{center}
  \end{columns}
  \begin{itemize}
    \item The objects being clustered can be customers, articles, images, genes, or users in a network.
    \item K-means is one way to turn those feature vectors into coherent groups.
  \end{itemize}
\end{frame}

\begin{frame}[t]{Representative Applications of Clustering}
  \begin{columns}[T,totalwidth=\textwidth]
    \column{0.48\textwidth}
      \begin{itemize}
        \item \textbf{Customer segmentation}: group customers by behavior, preferences, or demographics.
        \item \textbf{Document clustering}: organize large text corpora into coherent topics.
        \item \textbf{Anomaly detection}: identify unusual behavior that may indicate fraud or cyberattacks.
      \end{itemize}
    \column{0.48\textwidth}
      \begin{itemize}
        \item \textbf{Social network analysis}: find communities with similar links or interests.
        \item \textbf{Genomics}: discover genes with similar expression patterns.
        \item \textbf{Urban planning}: group regions with similar land use, demographics, or crime patterns.
      \end{itemize}
  \end{columns}

  \vspace{0.6em}
  \begin{block}{Why this matters}
    Clustering is useful whenever we want structure without labels: find groups, summarize populations, or flag unusual points.
  \end{block}
\end{frame}

\begin{frame}[t]{What K-Means Produces}
  \begin{columns}[T,totalwidth=\textwidth]
    \column{0.48\textwidth}
      \begin{center}
        \textbf{Before}

        \includegraphics[width=\linewidth,height=0.58\textheight,keepaspectratio]{kmeans_before.png}
      \end{center}
    \column{0.48\textwidth}
      \begin{center}
        \textbf{After}

        \includegraphics[width=\linewidth,height=0.58\textheight,keepaspectratio]{kmeans_after.png}
      \end{center}
  \end{columns}
  \vspace{0.25em}
  \begin{itemize}
    \item The algorithm assigns every point to exactly one cluster and summarizes each cluster by its center.
    \item The only thing it sees is distance, so the choice of features and scaling matters.
  \end{itemize}
\end{frame}

\begin{frame}[t]{The Optimization Problem}
  \begin{block}{K-means objective}
    \[
      W(\mathcal C,\mu_1,\ldots,\mu_k)
      =
      \sum_{\ell=1}^{k}\;\sum_{i:C(i)=\ell}\left\lVert x_i-\mu_\ell\right\rVert^2
    \]
  \end{block}

  \begin{itemize}
    \item $\mathcal C$ is the cluster assignment map and $\mu_\ell$ is the centroid of cluster $\ell$.
    \item Minimizing squared Euclidean distances makes clusters \textbf{compact} around their centers.
    \item For a fixed assignment, the best centroid is the sample mean:
  \end{itemize}

  \[
    \mu_\ell
    =
    \frac{1}{n_\ell}\sum_{j:C(j)=\ell} x_j,
    \qquad
    n_\ell = \#\{i : C(i)=\ell\}.
  \]
  \begin{itemize}
    \item This quantity is also called the \textbf{within-cluster scatter} or \textbf{WCSS}.
  \end{itemize}
\end{frame}

\begin{frame}[t]{Why the Mean Appears}
  \[
    \mu_\ell^\star
    =
    \arg\min_{m \in \mathbb R^d}
    \sum_{i:C(i)=\ell}\left\lVert x_i-m\right\rVert^2
    =
    \frac{1}{n_\ell}\sum_{i:C(i)=\ell}x_i
  \]

  \vspace{-0.25em}
  \begin{block}{Interpretation}
    With squared distance, the center that minimizes total within-cluster error is the arithmetic mean. That is why the algorithm is called \textbf{k-means}.
  \end{block}

  \begin{itemize}
    \item Assignment step: keep the centroids fixed and send each point to its nearest center.
    \item Update step: keep the assignments fixed and recompute each center as the mean of its assigned points.
    \item Each step decreases the objective, so the algorithm moves downhill until it reaches a local optimum.
  \end{itemize}
\end{frame}

\begin{frame}[t]{Why We Need an Iterative Heuristic}
  \begin{columns}[T,totalwidth=\textwidth]
    \column{0.53\textwidth}
      \begin{itemize}
        \item Searching over \emph{all} possible clusterings is combinatorial.
        \item The number of ways to partition $n$ points into $k$ nonempty clusters is the Stirling number:
      \end{itemize}
      \[
        S(n,k)=\frac{1}{k!}\sum_{\ell=1}^{k}(-1)^{k-\ell}\binom{k}{\ell}\ell^n
      \]
      \begin{itemize}
        \item For example, $S(10,4)=34105$ and $S(19,4)\approx 10^{10}$.
        \item Exhaustive optimization is unrealistic, so we use a fast greedy procedure instead.
      \end{itemize}
    \column{0.43\textwidth}
      \begin{block}{Lloyd's algorithm}
        \[
          C(i) \leftarrow \arg\min_{\ell}\lVert x_i-\mu_\ell\rVert^2
        \]
        \[
          \mu_\ell \leftarrow \frac{1}{n_\ell}\sum_{j:C(j)=\ell}x_j
        \]
      \end{block}

      \begin{itemize}
        \item Initialize $k$ centroids.
        \item Alternate assignment and update.
        \item Stop when assignments no longer change.
      \end{itemize}
  \end{columns}
\end{frame}

\begin{frame}[t]{A Concrete Run of Lloyd's Algorithm}
  \begin{columns}[T,totalwidth=\textwidth]
    \column{0.48\textwidth}
      \begin{center}
        \textbf{1. Choose \(k=4\)}

        \vspace{0.35em}
        \includegraphics[
          width=\linewidth,
          height=0.56\textheight,
          keepaspectratio
        ]{lloyd_choose_k.png}
      \end{center}
    \column{0.48\textwidth}
      \begin{center}
        \textbf{2. Randomly initialize centroids}

        \vspace{0.35em}
        \includegraphics[
          width=\linewidth,
          height=0.56\textheight,
          keepaspectratio
        ]{lloyd_random_init.png}
      \end{center}
  \end{columns}
  \begin{itemize}
    \item We start with a chosen value of \(k\) and an initial guess for the \(k\) centers.
    \item From here, the algorithm alternates between assignment and update.
  \end{itemize}
\end{frame}

\begin{frame}[t]{Iteration 1: Assign Every Point to Its Closest Center}
  \begin{center}
    \includegraphics[
      width=0.84\linewidth,
      height=0.68\textheight,
      keepaspectratio
    ]{lloyd_assign_step1.png}
  \end{center}
  \begin{itemize}
    \item Each point is assigned to the nearest current centroid.
    \item This creates a provisional partition of the data.
  \end{itemize}
\end{frame}

\begin{frame}[t]{Iteration 1: Recompute the Centroids}
  \begin{center}
    \includegraphics[
      width=0.84\linewidth,
      height=0.68\textheight,
      keepaspectratio
    ]{lloyd_update_step1.png}
  \end{center}
  \begin{itemize}
    \item Each centroid moves to the mean of the points currently assigned to it.
    \item In the plot, the red outlined markers show the newly computed centers.
  \end{itemize}
\end{frame}

\begin{frame}[t]{The Reassignment-Update Loop}
  \begin{columns}[T,totalwidth=\textwidth]
    \column{0.48\textwidth}
      \begin{center}
        \textbf{Another assignment step}

        \vspace{0.25em}
        \includegraphics[
          width=\linewidth,
          height=0.27\textheight,
          keepaspectratio
        ]{lloyd_assign_step2.png}

        \vspace{0.5em}
        \textbf{Near convergence}

        \vspace{0.25em}
        \includegraphics[
          width=\linewidth,
          height=0.27\textheight,
          keepaspectratio
        ]{lloyd_near_convergence.png}
      \end{center}
    \column{0.48\textwidth}
      \begin{center}
        \textbf{The following update}

        \vspace{0.25em}
        \includegraphics[
          width=\linewidth,
          height=0.27\textheight,
          keepaspectratio
        ]{lloyd_update_step2.png}

        \vspace{0.5em}
        \textbf{The final update}

        \vspace{0.25em}
        \includegraphics[
          width=\linewidth,
          height=0.27\textheight,
          keepaspectratio
        ]{lloyd_final_update.png}
      \end{center}
  \end{columns}
  \begin{itemize}
    \item Early on, assignments and centers can move a lot.
    \item As the algorithm settles, each update becomes smaller.
  \end{itemize}
\end{frame}

\begin{frame}[t]{Convergence}
  \begin{columns}[T,totalwidth=\textwidth]
    \column{0.48\textwidth}
      \begin{center}
        \textbf{Stop when the centers no longer move}

        \vspace{0.35em}
        \includegraphics[
          width=\linewidth,
          height=0.52\textheight,
          keepaspectratio
        ]{lloyd_final_update.png}
      \end{center}
    \column{0.48\textwidth}
      \begin{center}
        \textbf{Final cluster assignment}

        \vspace{0.35em}
        \includegraphics[
          width=\linewidth,
          height=0.52\textheight,
          keepaspectratio
        ]{lloyd_final_assignment.png}
      \end{center}
  \end{columns}
  \begin{itemize}
    \item Once the updated centers match the current centers, the algorithm has converged.
    \item The result is a locally optimal partition for this initialization.
  \end{itemize}
\end{frame}

\begin{frame}[t]{Initialization Matters}
  \begin{columns}[T,totalwidth=\textwidth]
    \column{0.38\textwidth}
      \begin{center}
        \textbf{Different random starts can lead to different answers}

        \vspace{0.5em}
        \includegraphics[width=0.98\linewidth,height=0.20\textheight,keepaspectratio]{init_different_start_1.png}

        \vspace{0.5em}
        \includegraphics[width=0.98\linewidth,height=0.20\textheight,keepaspectratio]{init_different_start_2.png}
      \end{center}
    \column{0.58\textwidth}
      \small
      \begin{itemize}
        \item K-means converges to a \textbf{local} optimum, so different starts can produce different partitions.
        \item A common fix is to run the algorithm many times and keep the best objective value.
        \item Intuition: if one region already has a center, we want the next center to land in a region that is still far away.
      \end{itemize}

      \begin{block}{k-means++ idea}
        \scriptsize
        Pick the first center uniformly. For each candidate point $x$, define
        $D(x)=\min_{j\in\{1,\ldots,i\}}\lVert x-m_j\rVert$,
        where $m_1,\ldots,m_i$ are the $i$ centers already chosen.

        So $D(x)$ is the distance from $x$ to its nearest chosen center, and the next center is sampled with
        $\Pr(x \text{ is chosen next}) \propto D^2(x)$.

        Points far from the current centers are therefore more likely to start a new cluster.
      \end{block}
  \end{columns}
\end{frame}

\begin{frame}[t]{Cluster Geometry}
  \begin{columns}[T,totalwidth=\textwidth]
    \column{0.46\textwidth}
      \includegraphics[width=\linewidth,height=0.48\textheight,keepaspectratio]{cluster_geometry_voronoi.png}
    \column{0.50\textwidth}
      \includegraphics[width=\linewidth,height=0.48\textheight,keepaspectratio]{cluster_geometry_nonconvex.png}
  \end{columns}
  \vspace{0.35em}
  \begin{itemize}
    \item With fixed centroids, k-means partitions the space into \textbf{Voronoi cells}: every point goes to its nearest center.
    \item Boundaries are piecewise linear, so clusters tend to be convex and roughly spherical in Euclidean space.
    \item That geometry makes elongated, imbalanced, and non-convex shapes harder to recover.
  \end{itemize}
\end{frame}

\begin{frame}[t]{When K-Means Works and When It Fails}
  \begin{columns}[T,totalwidth=\textwidth]
    \column{0.47\textwidth}
      \begin{block}{Works well when}
        \small
        \begin{itemize}
          \item clusters are compact and reasonably separated
          \item Euclidean distance is meaningful
          \item features are standardized and on comparable scales
          \item you want a fast baseline or preprocessing step
        \end{itemize}
      \end{block}
    \column{0.47\textwidth}
      \begin{block}{Struggles when}
        \small
        \begin{itemize}
          \item clusters are non-convex or have very different sizes
          \item there are strong outliers
          \item the true structure is density-based rather than centroid-based
          \item the value of $k$ is poorly chosen
        \end{itemize}
      \end{block}
  \end{columns}
  \vspace{0.15em}
  \begin{center}
    \includegraphics[width=0.68\linewidth,height=0.28\textheight,keepaspectratio]{kmeans_failure_example.jpg}
  \end{center}
\end{frame}

\begin{frame}[t]{Choosing \texorpdfstring{$k$}{k} with WCSS and the Elbow Method}
  \begin{columns}[T,totalwidth=\textwidth]
    \column{0.34\textwidth}
      \[
        \wcss
        =
        \sum_{j=1}^{K}\;\sum_{x_i\in C_j}\left\lVert x_i-\mu_j\right\rVert^2
      \]
      \small
      \begin{itemize}
        \item Smaller WCSS means tighter clusters.
        \item WCSS always decreases as $K$ increases.
        \item The elbow is where extra clusters give only a small additional improvement.
        \item The bend is a reasonable place to stop adding clusters.
      \end{itemize}
    \column{0.62\textwidth}
      \begin{center}
        \includegraphics[width=\linewidth,height=0.56\textheight,keepaspectratio]{elbow_method_plot.png}
      \end{center}
  \end{columns}
\end{frame}

\begin{frame}[t]{Choosing \texorpdfstring{$k$}{k} with Silhouette Scores}
  \begin{center}
    \begin{minipage}{0.80\textwidth}
      \footnotesize
      \[
        a(i)=\frac{1}{|C|-1}\sum_{x_j\in C,\;j\neq i} d(x_i,x_j)
      \]
      \[
        b(i)=\min_{C'\neq C}\frac{1}{|C'|}\sum_{x_j\in C'}d(x_i,x_j)
      \]
      \[
        s(i)=\frac{b(i)-a(i)}{\max\{a(i),b(i)\}},
        \qquad
        S=\frac{1}{n}\sum_{i=1}^{n}s(i)
      \]
      \begin{itemize}
        \item $a(i)$ measures cohesion inside the assigned cluster.
        \item $b(i)$ measures separation from the nearest competing cluster.
        \item $s(i)\approx 1$: the point is well matched to its own cluster.
        \item $s(i)\approx 0$: the point lies near a boundary between clusters.
        \item $s(i)<0$: the point may fit better in a different cluster.
      \end{itemize}
      \begin{center}
        Larger average silhouette means the clustering has both tighter groups and better separation between them.
      \end{center}
    \end{minipage}
  \end{center}
\end{frame}

\begin{frame}[t]{Using WCSS and Silhouette Together}
  \begin{columns}[T,totalwidth=\textwidth]
    \column{0.48\textwidth}
      \begin{block}{WCSS}
        \begin{itemize}
          \item Measures \textbf{compactness} inside each cluster.
          \item Great for elbow plots and fast screening of candidate $k$ values.
          \item Limitation: it ignores separation between different clusters.
        \end{itemize}
      \end{block}
    \column{0.48\textwidth}
      \begin{block}{Silhouette}
        \begin{itemize}
          \item Measures \textbf{compactness plus separation}.
          \item Great for comparing cluster quality more directly.
          \item Limitation: it is more computationally expensive.
        \end{itemize}
      \end{block}
  \end{columns}

  \begin{itemize}
    \item In practice, start with WCSS to narrow the range of plausible $k$ values.
    \item Then use silhouette plots, domain knowledge, and visual inspection to decide whether the clusters are actually meaningful.
  \end{itemize}
\end{frame}

\begin{frame}[t]{Interpreting the Clusters}
  \begin{columns}[T,totalwidth=\textwidth]
    \column{0.35\textwidth}
      \footnotesize
      \begin{itemize}
        \item K-means gives \textbf{partitions}, not semantic labels. Interpretation still requires subject-matter judgment.
        \item Inspecting cluster centroids often reveals the dominant profile of each group.
        \item In high dimensions, visualization often uses PCA, t-SNE, or UMAP after clustering.
        \item A good workflow is to standardize features, compare several $k$ values, inspect centroids and example members, then validate with domain knowledge.
      \end{itemize}
    \column{0.61\textwidth}
      \begin{center}
        \includegraphics[width=\linewidth,height=0.78\textheight,keepaspectratio]{domain_visual_categories.jpg}
      \end{center}
  \end{columns}
\end{frame}

\begin{frame}[t]{Takeaways}
  \begin{itemize}
    \item K-means is simple and fast because it alternates between assignment and least-squares centroid updates.
    \item Its objective is to minimize within-cluster squared distances to the centroids.
    \item Its Euclidean geometry favors compact, convex groups.
    \item Initialization and model selection matter, which is why k-means++ and validation tools are so important.
    \item Use it as a strong baseline, then move beyond it when the data violate its assumptions.
  \end{itemize}
\end{frame}

\end{document}
